\chapter{Time, the Trial, and the Study}
\label{chap:time-trial-study}

The first two chapters built numbers. This one builds the first quantity \emph{out of} numbers, and it
is the one every later measurement secretly depends on: time. The natural assumption is that an
instrument measures time the way a ruler measures a length --- that there is a flow out in the world
and the clock reads off how much of it has gone by. The construction takes the opposite view, and it
is worth stating plainly because it sets the tone for everything physical that follows. There is no
flow to read. A clock does not detect time; it \emph{makes} it, by imposing an irreversible order on
the facts it records. Time is the name we give to that imposed order, and a second is the name we give
to one of its steps.

The chapter follows the construction up four short rungs. First the clock complement, the minimal
device that manufactures a before and an after. Then the trial, a single repeatable measurement, and
the study, which pools many trials into a typical result. Then the way a study settles on an answer by
rounding to the nearest true, and what that does to our idea of cause. And finally the operational
second --- the point at which the construction stops describing time and simply defines it, a move we
will mark carefully as a modeling choice rather than a discovery.

\section{The clock complement}
\label{se:the-clock-complement}

The smallest thing that can keep time is a pair of opposites that take turns. The construction builds
exactly that and calls it a clock complement: two phase-opposite readings, a tick and a tock, a
before and an after, each carrying one of the two truth-values and each the negation of the other.
Nothing about this pair refers to an external duration. What it has instead is a step that takes the
current reading to the next one, and the step is built to be irreversible. A starting state can become
a response; a response cannot quietly revert to a starting state. The construction puts the point with
characteristic bluntness: a response cannot become an initial condition again without perjury. That
asymmetry --- forward is free, backward is forbidden --- is the whole of what the device knows about
time, and it is enough.

Reading such a clock is an act we already have a word for. To tell which phase the clock is in is to
\emph{tange} the two phases: to select, by the decision that holds of it, the one phase the clock is
currently presenting. And the question the construction actually asks of the pair --- do the before
and the after agree or disagree in truth-value? --- is itself decidable, a single yes-or-no about the
relative phase of the two readings. So time, at its root, is two opposed states, a decision that picks
out which one is showing, and a rule that the showing only ever advances. There is a quality here that
the construction cannot define and does not try to: what it is \emph{like} for one event to come
before another. It simply builds a device that recognizes the distinction reliably and gives it a
shared name, the way two people can agree to call something ``now'' without ever agreeing on what
nowness is.

A concrete run makes the manufacture of time visible. Begin a clock in its tick phase and step it:
tick gives way to tock, tock back to tick, and each step appends a record stamped only with the fact
that the new phase came after the old. Read those records back and they fall into a strict order ---
not because anything measured the gaps between them, but because the step refused to run backward, so
each record carries its predecessor and none can be slipped in earlier. Now start a second, identical
clock beside the first and let the two step at even slightly different rates. After a while they have
appended different numbers of records, and if you ask them how much time has passed they give
different answers. Neither is faulty. Each has been minting its own order rather than sampling a shared
one, and the gap between their counts is not an error to be calibrated away --- it is the simultaneity
of the two clocks coming apart, the whole of relativity's central surprise sitting in two flip-flops
that merely ran at their own pace.

This is precisely the lesson the construction files under Einstein's name. The Einstein experiment
makes a claim that sounds paradoxical until the clock complement makes it ordinary: temporal order
does not come from measuring an underlying flow; it comes from an instrument that enforces a directed
sequence of admissible records. An instrument produces time by imposing an irreversible ordering on
the facts it appends to its ledger, and for nothing else does it need to reach outside itself. Two
such instruments, run side by side, can disagree about how much time has passed precisely because each
is manufacturing its own order rather than both reading one shared flow --- which is the entire content
of relativity of simultaneity, recovered here from the bare fact that a clock is a thing that ticks
forward and not a window onto a river. The arrow of time, in this construction, is not assumed. It is
the irreversibility of a step, and that is something one can point at.

\section{The trial and the study}
\label{se:the-trial-and-the-study}

With a clock in hand, the construction can do the thing a clock is for: run the same measurement twice
and tell which came first. A single such measurement is a trial. A trial is a stimulus put to the
instrument and the settled response that comes back --- a question asked of the apparatus and the
answer it returns once it has stopped changing. The construction insists that a trial be repeatable,
which it makes precise by tying the stimulus and the response together into one inseparable
input-output pair: the question and its receipt are a single object, and a process that loses track of
which response went with which stimulus has stopped being a measurement.

One trial, though, settles almost nothing, and the construction is candid about why. A study is what
you get when you \emph{funge} many trials: pool the repetitions of the same measurement, collecting
the responses that share a stimulus into one bagged object, and ask what the pooled result typically
is. A study, notably, holds no numbers of its own --- it is a growing record of facts and trials, a
hypothesis with data accumulated behind it. Each new trial is appended; the study grows by funging one
more repetition into the pool. The typical response is not any single trial's answer but the shape the
pool takes as the trials accumulate.

The pooling does real work, and it is worth saying what. Each trial returns a settled response, but
each also carries a residue --- the leftover between what the apparatus reported and what the model
expected, in the precise sense the last chapter gave the word. The residues of repeated trials point
no consistent way: high here, low there, scattered around the thing being measured. Funge enough of
them and the scatter tends to cancel, while the genuine response, which points the same way every
time, survives the pooling intact. The typical response a study reports is what remains once the
per-trial residue has averaged itself out. That is why a study can name a sharper quantity than any of
its trials holds: it is not crowning a best trial but distilling the agreement across all of them, and
the sharpness it buys is paid for in the one currency the construction recognizes --- committed
events, appended until the scatter is small enough to ignore.

The reason a study is needed, rather than a single careful trial, is exactly the situation the
Dirichlet--Bancroft experiment dramatizes, and it is one the reader's phone performs every second. In
satellite positioning, the governing equations are exact and the satellite positions and relativistic
corrections are fully specified --- and yet three satellite receptions generally leave \emph{two}
geometrically valid solutions, one near the Earth's surface and one flung far into space. The
mathematics alone does not choose between them. Uniqueness arrives only when the ledger is
supplemented: another reception, or a standing constraint that the receiver is not in orbit. The
lesson the experiment draws is the one the study is built on. Model completeness and ledger
sufficiency are different things; even with exact equations, it is the \emph{accumulation of committed
events} that pins down a unique answer. A study funges trials for the same reason a receiver waits for
another satellite: a single observation underdetermines, and only the pool decides.

\section{Rounding to the nearest true}
\label{se:rounding-to-the-nearest-true}

How does a study, holding no numbers, ever deliver a verdict? It rounds. The order the construction
puts on studies prefers the richer, more-described hypothesis to the thinner one, and the practical
effect is that a study produces a rounded value of true: if the pooled trials make a claim seem true
enough, the study tends to accept it as true. This is an honest description of how measured conclusions
are actually reached --- not by attaining certainty, which the first chapter already ruled out for any
finite observer, but by accumulating enough agreeing trials that the nearest available verdict is
``true.'' Rounding to the nearest true is what a finite study does in place of proving.

The thing a study most wants to round its way to is a cause, and here the construction is careful, and
leans on Hume. A study tries to generalize a relation that holds independent of time --- does this
stimulus bring about that response? --- but Hume's warning is built into the design: experiencing the
response after the stimulus does not show that the one produces the other. The order is the order we
impose on the episodes, and what can actually be measured is not causation but how much truth there is
in the bare claim that the response follows the stimulus. The construction takes this seriously enough
that its studies never contain a primitive called cause; they contain only pooled trials and the
rounded verdict that one fact reliably follows another.

Why, then, does cause-and-effect feel so dependable, if it is never observed directly? The
Cause--Effect experiment gives the construction's answer, and it is a satisfying one. If events were
linked only by arbitrary correlation, the space of admissible futures would explode exponentially with
each refinement. They are not so linked: the requirement that the whole ledger stay globally consistent
collapses that space to a narrow band of predictable continuations. Reliable cause-and-effect, on this
reading, is not a law imposed from outside and not a miracle; it is the combinatorial residue of
admissibility itself --- what is left over once every inconsistent future has been struck out. A study
rounds toward causes because admissibility has already done most of the work of eliminating the
alternatives. The construction's name for running a study this way, as a process with inputs and
receipts that can be driven like a program, is a computation; and a study treated numerically, rather
than merely held as fact, is what lets the chapters ahead compute with measured time.

It is worth being precise about what that last step adds, because it is the bridge to every
calculation in the book. A study, by itself, is a verdict held as fact --- true, or true enough. It
becomes a computation the moment its questions are given inputs and its receipts are read as outputs:
the same pooled-trial machinery, now driven, a stimulus in and a settled response out, repeatable on
demand. And there is a careful distinction the construction draws here between being a number and being
usable as one. A number simply is what it is; a study made numeric does not pretend to be a number,
but it knows how it may be used as one --- which operations may be applied to its verdict without
overstating what the verdict will bear. That distinction is what keeps the chapters ahead honest when
they begin doing arithmetic on measured quantities: they compute with studies that carry their own
limits, not with bare numbers that have forgotten where they came from.

\section{The operational second}
\label{se:the-operational-second}

There is one move left, and it is the one to watch most closely, because it is where the construction
stops describing time and starts legislating it. Everything so far has given an order --- before,
after, irreversibly --- without giving a \emph{unit}. To get a second, the construction does the only
thing a finite device can: it picks a stable oscillation and declares one of its ticks to be the unit.
A second is one full swing of the clock complement, one tick-to-tock-to-tick of a flip-flop that has
been certified to repeat. Nothing in the world hands over this unit; it is manufactured by choosing an
oscillator and trusting it, and the trust is the substance of the definition.

That this is a definition and not a discovery is the heart of the Descartes experiment, and it deserves
to be stated without hedging. A fly crosses a continuous ceiling in continuous motion, but the
instrument that records the crossing does not admit the ceiling as such; it admits only pairs of
numbers read off calibrated scales, a quantized instance of position at each tick. Continuity survives
only as the limiting coherence of those discrete records. The experiment's verdict is exactly the
caution the operational second requires: quantization is an act of calibration, not a claim about the
world. The construction does not assert that time is discrete; it observes that measurement is
compelled to proceed discretely, and that quantization comes before analysis rather than after it. The
second is calibration made into a standard. It is as real and as useful as the oscillator it is pinned
to, and no more --- and a later instrument that adopts a different oscillator is keeping a different
second, which is a fact about clocks, not about time.

One consequence is worth drawing out, because the later chapters rely on it heavily. If a second is one
oscillator's certified tick --- a tange of that oscillator, with a funge of its repetitions standing
in for the unit --- then two instruments keep the same time only by agreeing to a convention that
lines their ticks up. There is no oscillator-independent fact that settles whether a tick here and a
tick there happened at once; simultaneity across a distance is not read off the world but established
by a rule, clocks exchanging records and adopting a procedure for calling their phases aligned. This
is the clock complement's opening lesson sharpened to its point: not only is the order of time
manufactured locally, one irreversible step at a time, so is the very claim that two separated events
coincide. The construction would rather buy that fact here, plainly, than smuggle it in later when it
begins comparing measurements made in different places.

\section{What is demonstrated, and what is assumed}
\label{se:demonstrated-assumed-ch3}

Following the discipline of the experimental library, we separate what the chapter built from what it
named. The temptation here is to think time has been explained; what has actually been built is
narrower and firmer than that.

What is demonstrated is a working clock and a working method. There is a clock complement: a pair of
phase-opposite states, a decidable reading that tanges the current phase, and an irreversible step
whose one-way character is a genuine property of the construction, not a stipulation. There is a trial,
a repeatable stimulus-and-settled-response, and a study that funges trials into a pooled verdict with a
total, computable order that prefers the better-supported hypothesis. Each of these is finite,
decidable, and checkable; none of it requires a primitive notion of duration, flow, or causation
anywhere.

\begin{quote}
\emph{A note on the labels.} Three names in this chapter carry more than the construction proves.
Calling the clock's imposed order \emph{time}, calling one tick of a chosen oscillator \emph{a second},
and calling a rounded study's verdict a statement of \emph{cause} are all modeling choices laid over
finite machinery. The order, the tick, and the rounding are earned; the identification of them with the
physicist's time, the metrologist's second, and the everyday relation of cause is interpretation. The
operational second in particular is honest only so long as we remember it is pinned to a specific
oscillator and travels no further than that oscillator's stability. The construction does not assume
the power to choose a unit from nowhere; it tanges a particular oscillator and funges its repetitions,
and the second is what that selection and pooling produce.
\end{quote}

What is assumed is the bridge from the device to the word. We assume that the irreversible order a
clock manufactures may be read as physical time, that a certified oscillator's tick may stand as the
universal second, and that a study's rounded verdict may be spoken of as a cause when it is really the
combinatorial residue of admissibility. These are the assumptions that let a finite ledger keep time
and reason about why things happen. With a clock built and trials in hand, the construction can finally
do what the opening of this chapter promised was coming: take rates of change seriously, and turn the
machinery of measurement on motion itself.
